\section{Related Work and the Novelty Test}
\label{sec:related}

\subsection{The Debt Family}
\label{sec:rel-family}
Technical debt began as a metaphor for expedient code~\cite{cunningham1992} and
matured into a research field with theory, taxonomy, and management
practice~\cite{kruchten2012,li2015,tom2013,alves2016,seaman2011,ernst2015}. The
family has since acquired members---architectural, test, documentation, security,
social, and hidden machine-learning debt~\cite{alves2016,rindell2019,tamburri2013,
sculley2015,avgeriou2016}---and, recently, organizational debt~\cite{alami2024};
practitioner literature gestures at ``operational debt''~\cite{owczarek2019} and,
closest in spirit to our subject, ``dark debt'': deferred, invisible operational
cost that reveals itself only in anomalies~\cite{stella2017}. Evidence debt gives
one species of dark debt a measurement theory. Evidence debt shares the family's defining structure
(present expedience, deferred cost, accruing carrying cost, principal/interest
decomposition) and differs in its \emph{object}: every existing debt type is a
property of the \emph{system or organization as it is}---its code, architecture,
tests, prose, security posture, structure---whereas evidence debt is a property of
the \emph{record of how it came to be}. The distinction has consequences the family's
tools cannot absorb: evidence debt is invisible to any inspection of the current
system state (a flawless codebase can carry catastrophic evidence debt); its
carrying cost is driven by third-party retention clocks and personnel turnover
rather than by ongoing development friction; and it uniquely exhibits mandatory
write-offs (\cref{prop:writeoff})---no amount of later diligence recovers a record
that no longer exists anywhere, whereas any code can, at some price, be refactored.

Self-admitted technical debt~\cite{potdar2014} offers an instructive contrast:
there, the record (a TODO comment) survives and the remediation is deferred; here,
the remediation (the system works) is complete and the \emph{record} is what was
deferred.

\subsection{Documentation Debt}
\label{sec:rel-doc}
The nearest family member deserves its own confrontation. Documentation debt names
missing or outdated \emph{descriptions of the system}: design docs, API references,
runbooks~\cite{alves2016}. Evidence debt names missing \emph{records of events}:
who changed what, when, why, under what authority, with what outcome.
The two differ in referent (state versus history), in falsifiability
(documentation can be regenerated from the system it describes; events cannot be
re-observed once their traces are gone), and in decay mechanics (documentation
rots by divergence from a living system~\cite{parnas1994,lehman1980}; evidence
rots by destruction of sources and loss of joinability). The software-maintenance
tradition has long measured the cost of understanding systems whose rationale is
lost~\cite{lientz1980,hunt2002}---``software archaeology'' is reconstruction
effort by another name, applied to code rather than operations. A team with
perfect current documentation and no change history has zero documentation debt
and maximal evidence debt.

\subsection{Observability and Logging}
\label{sec:rel-obs}
Observability practice~\cite{majors2022,otel,sigelman2010} and log
management~\cite{nist80092} concern the capture and interrogation of operational
telemetry, and are the mechanical substrate on which evidence capture rides;
empirical studies show that even this substrate is captured incompletely and
inconsistently in practice~\cite{yuan2012}. But
observability's design center is the \emph{present}: understanding the system's
current behavior, with retention windows sized for debugging, not governance. The
obvious objection---``is this not just missing logs?''---is answered by three
structural gaps.
First, logs are one store among six-plus; the deficits that dominated our
experiment---broken intent and authorization linkage, correlation rupture---live
\emph{between} stores, and no log completeness fixes joinability
(\cref{fnd:superadd}). Second, missing logs is a state description; evidence debt
is an economic quantity over a query workload with decay dynamics
(\cref{def:debt})---one can have complete logs and high debt (all short-retention)
or gappy logs and low debt (rare, low-stakes queries). Third, the observability
literature has no counterpart of irrecoverability, interest schedules, or
principal---the constructs that make deferral decisions analyzable. Secure-logging
research~\cite{schneier1999,crosby2009,haber1991} guarantees integrity of what is
kept and is orthogonal, as scoped in \cref{sec:intro}.

\subsection{Provenance, Forensics, and Records Management}
Provenance systems capture derivation graphs~\cite{provdm,muniswamy2006,
pasquier2017}---when deployed, they are evidence-debt \emph{prevention}
infrastructure, and our evidence graph (\cref{sec:model-graph}) is deliberately
PROV-compatible. What the provenance literature does not provide is the economics
of its own absence: what it costs to \emph{not} have deployed it, which is this
paper's subject. Digital forensics~\cite{casey2011,nist80086} supplies the
reconstruction craft and the source-exhaustion standard our irrecoverability
certification borrows; its setting is adversarial and case-bound, whereas evidence
debt is an estate-level, mostly non-adversarial accounting. Records management and
archival science~\cite{iso15489,oais,rothenberg1995} have long managed retention
against future obligation; evidence debt can be read as bringing their concern into
software operations with a formal cost model attached, where change velocity and
fragmentation are orders of magnitude beyond the archival setting.

\subsection{Knowledge Management and Turnover}
Organizational memory research established that organizations forget through
personnel flow~\cite{walsh1991,nonaka1994}, and empirical software engineering has
quantified turnover-induced knowledge loss~\cite{rigby2016} and truck
factors~\cite{avelino2016}. In our model these appear precisely as the hazard
survival of human recovery paths (\cref{sec:model-decay})---the term that makes
interest smooth between cliffs (\cref{fig:interest}). The debt frame adds what the
knowledge literature lacks for this domain: a per-change, schema-relative,
countable unit of deferral (the deficit) connecting individual operational
decisions to the aggregate forgetting rate.

\subsection{Linkage, Traceability, Process Mining, and Data Quality}
\label{sec:rel-linkage}
Our reconstruction theory imports record linkage~\cite{fellegi1969,christen2012}
(effort/error structure of inferential joins), the missing-data
tradition~\cite{rubin1976,little2002} (missingness as a mechanism, not just a
rate---operational evidence is emphatically not missing at random: P1--P8 delete
specific interrogatives under specific pressures, a bias field studies must model),
and data-quality theory~\cite{wang1996} (fitness-for-use relativity, mirrored in
our schema- and workload-relativity).

Three empirical traditions have, in fact, already measured instances of the
phenomenon this paper names, and confronting them sharpens both the novelty claim
and its limits. The mining-software-repositories community documented
\emph{missing links} between bugs and fixing commits---our P5, correlation
rupture, in the wild---and showed that heuristically reconstructed link sets are
not merely incomplete but systematically \emph{biased}, corrupting downstream
conclusions~\cite{bachmann2010,bird2009}: field evidence that the false-join
channel of \cref{fnd:false} operates outside simulations. Requirements-%
traceability research has treated missing artifact links and their after-the-fact
recovery for three decades~\cite{gotel1994,antoniol2002}, with the same
effort/precision economics. Process mining reconstructs operational processes
from partial, noisy event logs and maintains an explicit log-quality
literature~\cite{aalst2016}. What none of these supplies---and what evidence debt
adds over them---is the estate-level economic aggregation: a schema-relative
deficit census, decay dynamics with principal/interest structure, an
irrecoverability boundary, and a workload-weighted cost model connecting today's
capture decisions to tomorrow's reconstruction burden. Conversely, these
literatures ground our claim that the phenomenon is real and measurable at
scale; they are the natural execution venues for parts of \cref{sec:field}.

Configuration drift and IaC research~\cite{rahman2019mapping,beetz2022,
opengitops} concern divergence of running state from desired state; evidence debt
concerns divergence of \emph{recorded history} from \emph{actual history}---a
system with zero configuration drift can be maximally evidence-indebted about how
it got there. Insecure IaC practice itself illustrates deficit creation at the
source~\cite{rahman2019}. Supply-chain integrity work~\cite{torres2019,slsa}
makes particular evidence chains cryptographically strong; it hardens what is
captured rather than pricing what is not. Regulatory control catalogs demand
auditable change records outright~\cite{nist80053,sox2002}; evidence debt supplies
the quantity those obligations implicitly assume is being managed.

\subsection{The Novelty Test, Answered Directly}
\label{sec:novelty}

\emph{Why is this not simply technical debt?} Different object (record of history
versus state of system), different decay drivers (third-party clocks and turnover
versus development friction), different terminal behavior (mandatory write-off
versus indefinite refactorability), different observability (invisible to system
inspection). \emph{Why not documentation debt?} State versus events;
regenerability versus unrepeatable observation (\cref{sec:rel-doc}). \emph{Why
not missing logs?} Logs are one store; the phenomenon is cross-store joinability,
workload-relative cost, and decay dynamics; completeness of capture at one point
neither implies nor is implied by low debt (\cref{sec:rel-obs}). \emph{Why not
observability debt?} Insufficient observability is one \emph{creation condition}
among the eight of \cref{sec:practices}; the debt also grows from retention,
migration, turnover, and identifier rupture after perfectly observable capture.
\emph{What unique phenomenon does evidence debt capture?} The deferred,
decaying, workload-contingent cost of reconstructing operational history---a
quantity with its own creation conditions (RQ1), estimator theory (RQ2),
effort-before-failure dynamics (RQ3, \cref{fnd:effort}), and irrecoverability
boundary (RQ4, \cref{fnd:redundancy,fnd:constructs})---none of which is expressible as a
property of the current system state that the existing debt family describes.

\begin{figure}[t]
\centering
\resizebox{\columnwidth}{!}{% Technical debt vs evidence debt (single column, resized)
\begin{tikzpicture}[node distance=2.8mm and 4mm]
  % shared structure
  \node[lbl, anchor=west] (h0) at (0,0) {\textbf{Shared deferral structure}};
  \node[stage, below=2.5mm of h0.west, anchor=north west, minimum width=20mm] (a)
    {expedient act\\{\smlbl now}};
  \node[stage, right=4mm of a, minimum width=20mm] (b)
    {carrying cost\\{\smlbl accrues}};
  \node[stage, right=4mm of b, minimum width=20mm] (c)
    {repayment event\\{\smlbl later, contingent}};
  \draw[flow] (a) -- (b);
  \draw[flow] (b) -- (c);

  % divergence table
  \matrix (m) [below=5mm of b, matrix of nodes, column sep=2.5mm, row sep=0.7mm,
    nodes={font=\fontsize{6.8}{7.9}\selectfont, anchor=west, align=left},
    column 1/.style={nodes={font=\fontsize{6.8}{7.9}\selectfont\bfseries}}]
  {
    \strut & \textbf{Technical debt} & \textbf{Evidence debt} \\
    object & system as it is (code, design) & record of how it came to be \\
    visible via & inspection of current system & only via queries against history \\
    interest driver & ongoing development friction & retention clocks, turnover, migration \\
    growth shape & roughly continuous & plateaus and cliffs (Prop.~3) \\
    terminal evidence state & typically remediable while artifacts remain & physical irrecoverability (Prop.~4) \\
    policy-burden channels & engineering effort & effort, accepted error, abstention \\
  };
  \draw[black!50] ($(m-1-1.north west)+(-1mm,0.4mm)$) -- ($(m-1-3.north east)+(1mm,0.4mm)$);
  \draw[black!50] ($(m-1-1.south west)+(-1mm,-0.2mm)$) -- ($(m-1-3.south east)+(1mm,-0.2mm)$);
  \draw[black!50] ($(m-7-1.south west)+(-1mm,-0.6mm)$) -- ($(m-7-3.south east)+(1mm,-0.6mm)$);
\end{tikzpicture}
}
\caption{Technical debt versus evidence debt. Both share the deferral structure
(left): an expedient act now, a carrying cost, a repayment event. They diverge in
every load-bearing property (right): object, visibility, interest driver, growth
shape, terminal state, and repayment currency. The comparison is the condensed form
of \cref{sec:novelty}.}
\label{fig:tdvsed}
\end{figure}

\Cref{fig:tdvsed} condenses the comparison. The claim, precisely bounded: evidence
debt is a member of the debt family by structure, and reducible to no existing
member by object, dynamics, or management levers.
